\documentclass{article}
\usepackage[portuguese]{babel}
\usepackage{listings}
\usepackage{xcolor}

\definecolor{codegreen}{rgb}{0,0.6,0}
\definecolor{codegray}{rgb}{0.5,0.5,0.5}
\definecolor{codepurple}{rgb}{0.8,0,0.2}
\definecolor{backcolour}{rgb}{.8,.8,1}

\lstdefinestyle{mystyle}{
    backgroundcolor=\color{backcolour},   
    commentstyle=\color{codegreen},
    keywordstyle=\color{blue},
    numberstyle=\tiny\color{codegray},
    stringstyle=\color{codepurple},
    basicstyle=\ttfamily\footnotesize,
    breakatwhitespace=false,         
    breaklines=true,                 
    captionpos=b,                    
    keepspaces=true,                 
    numbers=left,                    
    numbersep=5pt,                  
    showspaces=false,                
    showstringspaces=false,
    showtabs=false,                  
    tabsize=2
}

\lstset{style=mystyle}


\begin{document}

\section{Estrutura de repeticao}

Para que o nossos algoritmos fiquem mais faceis de testar. Para que nao tenhamos que executar o arquivo repetidas vezes, vamos criar uma estrutura para repetir nossos algoritmos ate que desejamos parar.

Para isso vamos criar uma variavel para indicar se o usuario deseja parar

    \begin{lstlisting}[language=Python]
    deveParar = ""\end{lstlisting}

Agora criamos um laco while para repetir nosso algoritmo.

\begin{lstlisting}[language=Python]
    while(deveParar == ""):
        #Algoritmo a ser executado.
        
        deveParar = input("deseja continuar?(enter para continuar, qualquer tecla para parar)\n")
        
    print("Ate mais!")\end{lstlisting}

Observe que quando o usuario apenas teclar enter, a funcao \textit{input()} vai retornar um texto vazio ````, assim nosso algoritmo sera repetido, mas quando qualquer outra tecla for pressionada, a expressao \textit{[(deveRepetir == ````) False]}, dessa forma sairemos do laco de repeticao e o programa ira continuar, ou no nosso caso parar, pois nao tem mais nada depois do laco de repeticao

\end{document}

